\par\noindent\textbf{(1)} Euler式 \(I_i\dot\omega_i=(I_j-I_k)\omega_j\omega_k\) 对 \((i,j,k)\) 循环成立。乘以 \(\omega_i\) 求和得 \(\dot E=\frac12d(\sum I_i\omega_i^2)/dt=0\)，乘以 \(I_i\omega_i\) 求和得 \(\frac12d(\sum I_i^2\omega_i^2)/dt=0\)。这分别是能量椭球与动量球的交线约束。
\par\noindent\textbf{(2)} 在体轴角动量分量 \(L_i=I_i\omega_i\) 中，两曲面为 \(\sum L_i^2=L^2\) 和 \(\sum L_i^2/I_i=2E\)。惯量椭球切平面的法线沿 \((L_1/I_1,L_2/I_2,L_3/I_3)=\boldsymbol\omega\)，从而瞬时角速度与接触几何相关；实际轨迹为两曲面交线，不可理解为两曲面随意相交。
\par\noindent\textbf{(3)} 在 \(\omega_2=\Omega,\omega_1=u,\omega_3=v\) 附近，线性化得 \(\dot u=(I_2-I_3)\Omega v/I_1\)、\(\dot v=(I_1-I_2)\Omega u/I_3\)。故 \(\ddot u=\lambda^2u\)，\(\lambda^2=\Omega^2(I_3-I_2)(I_2-I_1)/(I_1I_3)>0\)，表现指数不稳定。
\par\noindent\textbf{(4)} 绕 \(I_1\) 轴线性频率平方为 \(\Omega^2(I_3-I_1)(I_2-I_1)/(I_2I_3)\)，绕 \(I_3\) 轴为 \(\Omega^2(I_3-I_1)(I_3-I_2)/(I_1I_2)\)，均正。若两惯量趋于相同，其中一频率趋零，意味着轴对称简并而不是出现固定非零频率。
